\subsubsection{面向小 Batch 场景的量化路径优化}

为提高量化稀疏方法在小 batch 场景下的竞争力，本文在 `input=4096, output=256` 的设置下，对 2bit 量化路径进行了定向运行时调优，并使用隔离式 batch-size runner 对 `BS=1..8` 的五组配置重新测量吞吐量。结果表明，量化路径的优势主要集中在 `BS=2..6` 区间：`BS=2` 时最优配置为 `Sparse50+2bit`，吞吐量为 `42.76` tokens/s；`BS=3..6` 时最优配置均为 `Sparse70+2bit`，吞吐量分别达到 `57.81`、`66.42`、`74.22` 和 `80.38` tokens/s。

不过，这一优势并未在全部 batch size 上持续扩大。对于 `Sparse50` 路径，`Sparse50+2bit` 只在 `BS=1..5` 上保持领先，在 `BS=6` 后被 sparse-only 版本追平并反超；对于 `Sparse70` 路径，`Sparse70+2bit` 在 `BS=1..7` 仍优于 `Sparse70`，但在 `BS=8` 时后者以 `96.43` tokens/s 重新取得最优。因此，这组结果更准确地说明：当前优化后的量化路径已经能够在小到中等 batch 区间稳定释放端到端收益，但在更大 batch 下仍需进一步压缩量化侧自身的附加开销。
